\begin{figure}[!htbp]
  \centering
  \begin{tikzpicture}[
    every node/.style={font=\small, align=center},
    head/.style={rectangle, rounded corners=4pt, draw=red!60, very thick,
                 fill=red!8, minimum width=11.2cm, minimum height=0.95cm},
    col/.style={rectangle, rounded corners=3pt, draw=blue!55!black, thick,
                fill=blue!8, minimum width=3.2cm, minimum height=3.2cm},
    foot/.style={rectangle, rounded corners=3pt, draw=gray!65, thick,
                 fill=gray!10, minimum width=11.2cm, minimum height=1.0cm},
    arr/.style={-stealth, thick}
  ]
    \node[head] (head) at (0,3.6) {\textbf{同一 behavior-guided 校准框架在不同 bit-width 约束下的路径实例化}};

    \node[col] (int8) at (-4.0,1.0) {\textbf{INT8 规范路径}\\[3pt]
      对称 per-group\\
      静态 scale\\
      自适应保护\\[3pt]
      \footnotesize 规范验证实例};
    \node[col] (int4s) at (0,1.0) {\textbf{对称 \texttt{INT4} 试探路径}\\[3pt]
      对称 per-group\\
      位宽从 8 降到 4\\
      Key 排序更脆弱\\[3pt]
      \footnotesize 暴露格式局限};
    \node[col] (role) at (4.0,1.0) {\textbf{\texttt{INT4-RoleAlign}}\\[3pt]
      K: per-channel\\
      V: per-token\\
      K/V split proxy\\[3pt]
      \footnotesize role-aware low-bit path};

    \node[foot] (foot) at (0,-1.9) {\textbf{共享离线校准工作流与稳健选择纪律}\\
      \footnotesize 变的是路径接口、量化轴与行为代理；不变的是 offline workflow、feasible-set 约束与 artifact discipline};

    \draw[arr] (head.south west) |- (int8.north);
    \draw[arr] (head.south) -- (int4s.north);
    \draw[arr] (head.south east) |- (role.north);
    \draw[arr] (int8.south) -- (foot.north west);
    \draw[arr] (int4s.south) -- (foot.north);
    \draw[arr] (role.south) -- (foot.north east);
  \end{tikzpicture}
  \caption{跨位宽路径的实例化关系与统一校准工作流。图中强调：INT8 规范路径、对称 \texttt{INT4} 试探路径与 \texttt{INT4-RoleAlign} 共享同一离线校准 workflow 与稳健选择纪律，但在路径接口、量化轴与行为代理上逐步分化。}
  \label{fig:ch3-path-instantiation}
\end{figure}
